AMP addresses this problem through Automated Meter Profiling with Latent Distribution Profiling (AMP-LDP). Each customer history is divided into 30-day measurement windows. Every window is represented by 17 frozen summary features and passed through a learned $17\!\rightarrow\!32\!\rightarrow\!16$ encoder. For each latent dimension, AMP-LDP retains the mean, stabilized population standard deviation, 10th percentile, median, 90th percentile, and maximum across the customer's windows. The resulting 96-dimensional profile is classified with histogram gradient boosting, and the primary customer probability is the arithmetic mean of five frozen seed pipelines. Supervision is applied at customer level; no instance-level theft loss is imposed.

The study separates two questions that are easily conflated. First, does preserving a distribution of window behaviour avoid the information bottleneck observed with single-vector pooling? Second, after the aggregation rule is held fixed, does a learned latent representation add information beyond applying the same six statistics directly to the 17 handcrafted window summaries? The DIRECT control answers the second question, while the engineered FULL\_48 customer profile provides a strong established comparator. The completed evaluation uses frozen development, confirmation, and final-evaluation roles rather than selecting a model from the terminal cohort.

The primary research question is:

\begin{quote}
Under customer-level SGCC theft labels and no verified theft-event timestamps, does learning representations of unlabeled consumption windows before distribution-preserving customer aggregation improve customer-level discrimination beyond applying the same aggregation directly to handcrafted window features?
\end{quote}

The report makes three bounded contributions. First, it formulates electricity-theft screening at the granularity supported by the data: labelled customers containing unlabeled temporal windows. Second, it investigates a distribution-preserving learned customer representation and evaluates it against both a matched DIRECT aggregation control and a strong engineered FULL\_48 comparator. Third, it reports a governed evidence chain across development, confirmation, and final evaluation using customer-level ROC-AUC, Average Precision, paired uncertainty, and seed-robustness checks. These contributions concern customer-level discrimination on the governed SGCC cohorts; they do not establish theft-period localization, window-level theft probabilities, causal mechanisms, external-utility generalization, or universal state-of-the-art performance.
